% VI32-DETAILED-PROBLEM-09
\begin{problem}[VI.32.P9: Crossing axes]\label{prob:vi32-09}
Analyze regularity at the origin of
\[
V(xy)\subset\mathbb A^2_k.
\]

\textbf{Provenance.} Canonical VI/32 synthesis.
\end{problem}

\ifdefined\FullProblemDossiers
\begin{solution}
The coordinate ring is
\[
A=k[x,y]/(xy).
\]
Its minimal primes \((x)\) and \((y)\) correspond to the two coordinate axes. Each component has dimension \(1\), so
the local ring at the origin has Krull dimension \(1\).

For
\[
f=xy,
\]
the gradient is
\[
\nabla f=(y,x),
\]
which vanishes at the origin. Therefore
\[
T_0X=k^2,
\qquad
\dim_kT_0X=2.
\]
Thus
\[
1=\dim\mathcal O_{X,0}<2=\dim_kT_0X,
\]
and the origin is singular.

The singularity mechanism differs from the cusp: here the scheme is reducible and two regular one-dimensional
components meet. The initial form is already
\[
xy,
\]
so the tangent cone is the union of the two coordinate axes themselves.
\end{solution}
\fi
